\chapter{标架、余标架与非完整性}

切空间在每一点都是线性空间，但实际计算必须选择局部基。坐标基 $\{\partial_i\}$ 只是一类特殊标架：它们来自坐标函数，因此彼此对易。许多对称性明显的计算反而更适合使用正交、左不变或随体标架，这些基一般不来自任何坐标系。本章只研究标架自身的运动学结构，特别是 Lie 括号和余标架外微分如何编码“非坐标性”；联络仍留到下一章。

\section{局部标架与对偶余标架}
在开集 $U\subset M$ 上，一组处处线性无关的向量场
\[
\{e_a\}_{a=1}^n
\]
称为局部标架。它给出每点切空间的一组基。唯一的对偶一形式 $\{\theta^a\}_{a=1}^n$ 满足
\begin{equation}\label{eq:math-dg-frame-duality}
\theta^a(e_b)=\delta^a{}_b.
\end{equation}
任意向量场和一形式分别展开为
\[
X=X^ae_a,
\qquad
\alpha=\alpha_a\theta^a.
\]
指标只是相对于所选标架的分量标签；在尚未引入度量时，上、下指标之间没有“升降”关系。

若从标架 $e_a$ 换到另一标架 $e'_a$，用可逆矩阵值函数 $\Lambda=(\Lambda^b{}_a)$ 写成
\begin{equation}\label{eq:math-dg-frame-change}
e'_a=e_b\Lambda^b{}_a.
\end{equation}
对偶条件\eqnrefs{eq:math-dg-frame-duality}要求
\[
\theta'^a=(\Lambda^{-1})^a{}_b\theta^b.
\]
因此同一个向量的分量满足
\[
X'^a=(\Lambda^{-1})^a{}_bX^b,
\]
而同一个一形式的分量满足
\[
\alpha'_a=\alpha_b\Lambda^b{}_a.
\]
几何对象不随换基改变，改变的是表示它的分量。


更一般地，一个 主 $G$-丛 是带自由右 $G$-作用的纤维丛
\[
\pi:P\to M,
\]
其每根纤维都是一个 $G$-扭子空间：不存在优先选定的单位点，但任意两点由唯一群元素联系。局部平凡化给出
\[
\pi^{-1}(U_\alpha)\cong U_\alpha\times G.
\]
在重叠区域 $U_\alpha\cap U_\beta$ 上，两种局部截面满足
\begin{equation*}s_\beta=s_\alpha g_{\alpha\beta},
\qquad
g_{\alpha\beta}:U_\alpha\cap U_\beta\to G,\end{equation*}
并在三重重叠上满足 上闭链条件
\begin{equation}
g_{\alpha\beta}g_{\beta\gamma}=g_{\alpha\gamma}.
\label{eq:principal-transition-cocycle}
\end{equation}
因此 主丛 的全局拓扑被局部平凡化之间的粘合函数编码。

给定群表示
\[
\rho:G\to GL(W),
\]
可以由 主丛 构造 伴随 向量丛
\begin{equation*}E=P\times_\rho W
:=(P\times W)/\sim,
\qquad
(pg,w)\sim(p,\rho(g)w).\end{equation*}
标架丛是母对象：取 $P=\operatorname{Fr}(M)$、$G=GL(n,\RR)$，标准表示给回切丛
\[
TM\cong\operatorname{Fr}(M)\times_{GL(n)}\RR^n,
\]
对偶表示给出 $T^*M$，张量表示给出任意 $T^r{}_sM$。因此“张量分量怎样随标架变换”本质上只是 伴随-丛 表示律。

\begin{derivation}[局部截面、转移函数与伴随丛分量]
主丛局部截面之间的换页函数完全由结构群的右作用决定；伴随丛 上截面的分量变换则由表示 $\rho$ 把这一右作用翻译成 $W$ 上的线性变换。两步合起来给出"换标架"的统一机制。

设 $s_\alpha:U_\alpha\to P$、$s_\beta:U_\beta\to P$ 是开覆盖 $\{U_\alpha\}$ 上的两个局部截面。在重叠 $U_\alpha\cap U_\beta$ 上，$s_\alpha(p)$ 与 $s_\beta(p)$ 落在同一根 纤维 $P_p$ 中。由于结构群 $G$ 在每根 纤维 上的右作用自由且传递，对每个 $p\in U_\alpha\cap U_\beta$ 存在唯一 $g_{\alpha\beta}(p)\in G$ 使
\begin{equation*}
 s_\beta(p)=s_\alpha(p)\,g_{\alpha\beta}(p),
 \qquad g_{\alpha\beta}:U_\alpha\cap U_\beta\to G.
\end{equation*}
光滑性由局部截面的光滑性与右作用的光滑性保证。$g_{\alpha\beta}$ 称为 \emph{转移函数}（或换页函数）。

在三重重叠 $U_\alpha\cap U_\beta\cap U_\gamma$ 上依次换页，
\begin{equation*}
 s_\gamma
 =s_\beta\,g_{\beta\gamma}
 =\bigl(s_\alpha\,g_{\alpha\beta}\bigr)g_{\beta\gamma}
 =s_\alpha\,\bigl(g_{\alpha\beta}g_{\beta\gamma}\bigr).
\end{equation*}
另一方面直接换页给出 $s_\gamma=s_\alpha\,g_{\alpha\gamma}$。由右作用的自由性（若 $s_\alpha g=s_\alpha h$ 则 $g=h$），得到
\begin{equation*}
 g_{\alpha\gamma}=g_{\alpha\beta}\,g_{\beta\gamma},
\end{equation*}
即\eqnrefs{eq:principal-transition-cocycle}。由 $s_\alpha=s_\alpha g_{\alpha\alpha}$ 及右作用自由性得 $g_{\alpha\alpha}=e$；把 $s_\beta=s_\alpha g_{\alpha\beta}$ 与 $s_\alpha=s_\beta g_{\beta\alpha}$ 合并再用自由性得 $g_{\alpha\beta}=g_{\beta\alpha}^{-1}$。$\{g_{\alpha\beta}\}$ 因此是取值于 $G$ 的 \v Cech $1$-上闭链，主丛的拓扑分类恰由其 \v Cech 上同调类决定。

设 $E=P\times_\rho W$，$\sigma$ 是 $E$ 上一个截面。在 $U_\alpha$ 上把它写成 $\sigma|_{U_\alpha}=[s_\alpha,w_\alpha]$，其中 $w_\alpha:U_\alpha\to W$ 是局部 $W$-值函数。在重叠上同一截面有两种表示
\begin{equation*}
 [s_\alpha(p),w_\alpha(p)]=[s_\beta(p),w_\beta(p)].
\end{equation*}
把 $s_\beta=s_\alpha g_{\alpha\beta}$ 代入右端，再用等价关系的定义 $(p g,w)\sim(p,\rho(g)w)$ 把 $g_{\alpha\beta}$ 从第一分量搬到第二分量：
\begin{equation*}
 [s_\alpha,w_\alpha]
 =[s_\alpha g_{\alpha\beta},w_\beta]
 =[s_\alpha,\rho(g_{\alpha\beta})\,w_\beta].
\end{equation*}
由 $P\times W\to E$ 在每个 $p$ 处限制为 $W\to E_p$ 的单射性，得到分量变换律
\begin{equation*}
 w_\alpha=\rho(g_{\alpha\beta})\,w_\beta,
 \qquad
 w_\beta=\rho(g_{\alpha\beta}^{-1})\,w_\alpha
 =\rho(g_{\beta\alpha})\,w_\alpha.
\end{equation*}

取 $P=\operatorname{Fr}(M)$、$G=GL(n,\RR)$：
\begin{itemize}
\item 标准表示 $\rho=\mathrm{id}$ 给出 $w_\alpha=g_{\alpha\beta}w_\beta$，即切向分量按 $GL(n,\RR)$ 矩阵换标架，对应 $E=TM$。
\item 对偶表示 $\rho(g)=(g^{-1})^T$ 给出 $w_\alpha=(g_{\alpha\beta}^{-1})^T w_\beta$，对应 $E=T^*M$，余切分量按逆的转置变换。
\item $(r,s)$ 张量表示给出 $w_\alpha=g_{\alpha\beta}^{\otimes r}\otimes(g_{\alpha\beta}^{-1})^{\otimes s}w_\beta$，对应 $E=T^r{}_sM$。
\end{itemize}
因此所有换标架分量公式——向量、余向量、张量、张量密度——都是同一条 伴随-丛 表示律 $w_\alpha=\rho(g_{\alpha\beta})w_\beta$ 在不同表示 $\rho$ 下的特例。

规范理论中物质场是伴随丛 $P\times_\rho W$ 的截面，$W$ 是内部空间（如 $\CC^n$ 上的表示空间），$g_{\alpha\beta}$ 是规范变换。分量变换律 $w_\alpha=\rho(g_{\alpha\beta})w_\beta$ 正是"波函数在重叠区域上的规范变换"的几何来源。
\end{derivation}

局部标架还可以提升为一个全局几何对象。定义标架丛
\[
\operatorname{Fr}(M):=\bigsqcup_{p\in M}
\operatorname{Iso}(\RR^n,T_pM),
\]
其一点 $u\in\operatorname{Fr}(M)$ 是一个线性同构 $u:\RR^n\to T_pM$，也就是 $p$ 处的一组有序基。投影 $\pi(u)=p$，右作用
\[
R_A(u):=u\circ A,
\qquad A\in GL(n,\RR),
\]
使 $\operatorname{Fr}(M)$ 成为主 $GL(n,\RR)$ 丛。一个局部标架 $e_a$ 与一个局部截面
\[
s:U\to\operatorname{Fr}(M),
\qquad
s(p)(\mathbf e_a)=e_a(p)
\]
完全等价，因此“换标架”就是主丛截面之间的 规范 变换，而不是坐标变换的同义词。

标架丛上存在不需要任何联络就能定义的典范 $\RR^n$-值一形式，称为 solder 形式 或 tautological 形式：
\begin{equation*}\vartheta_u(V)
:=u^{-1}(\pi_*V),
\qquad
V\in T_u\operatorname{Fr}(M).\end{equation*}
它把标架丛上的切向量投到底流形，再用当前标架 $u$ 读成 $\RR^n$ 分量。若 $s$ 是由 $e_a$ 给出的局部截面，则
\begin{equation*}s^*\vartheta
 =\theta^a\mathbf e_a.\end{equation*}
所以普通余标架并不是凭空出现的局部记号，而是全局 solder 形式 在局部 规范 中的表示。由右作用定义还可直接得到
\begin{equation*}R_A^*\vartheta=A^{-1}\vartheta.\end{equation*}
下一章的联络一形式将与 $\vartheta$ 一起组成 Cartan 结构方程；前者描述“如何比较邻点标架”，后者描述“标架如何焊接到底流形切空间”。

\section{结构函数与余标架外微分}
一般标架不必对易。定义结构函数 $C^c{}_{ab}\in C^\infty(U)$ 为
\begin{equation}\label{eq:math-dg-synthesis-anholonomy}
[e_a,e_b]=C^c{}_{ab}e_c,
\qquad
C^c{}_{ab}=-C^c{}_{ba}.
\end{equation}
若 $e_a=\partial/\partial y^a$ 来自某组局部坐标，则偏导数交换，故全部 $C^c{}_{ab}=0$。反过来，一组线性无关且彼此对易的向量场局部上可以同时化成坐标基；严格的可积性机制将在 Frobenius 定理中给出。

结构函数与余标架外微分包含完全相同的信息。

\begin{derivation}[从 Lie 括号得到 \eqnrefs{eq:math-dg-frame-structure-equation} 中的 $\dd\theta^a$]
利用外微分在两个向量场上的内在公式，将 $\dd\theta^a(e_b,e_c)$ 用 Lie 括号 $[e_b,e_c]$ 表达，再用结构函数 $C^a{}_{bc}$ 代入。

对任意一形式 $\alpha$，外微分在两个向量场上满足
\begin{equation*}
 \dd\alpha(X,Y)
 =X[\alpha(Y)]-Y[\alpha(X)]-\alpha([X,Y]).
\end{equation*}
这是外微分的无坐标定义，可直接验证：在坐标基下两端都等于 $\partial_i\alpha_j-\partial_j\alpha_i$，由张量性全局成立。第三项中 Lie 括号的出现正是外微分"反对称化方向导数并扣除方向场自身不交换"的几何实质。

令 $\alpha=\theta^a$、$X=e_b$、$Y=e_c$。由余标架的对偶性 $\theta^a(e_b)=\delta^a{}_b$（\eqnrefs{eq:math-dg-frame-duality}），
\begin{equation*}
 e_b[\theta^a(e_c)]
 =e_b(\delta^a{}_c)=0,
 \qquad
 e_c[\theta^a(e_b)]
 =e_c(\delta^a{}_b)=0,
\end{equation*}
因为 $\delta^a{}_c$ 是常数，其方向导数为零。故前两项消失，只剩
\begin{equation*}
 \dd\theta^a(e_b,e_c)
 =-\theta^a([e_b,e_c]).
\end{equation*}
这一步只用对偶性，不需要任何联络或度量。

由\eqnrefs{eq:math-dg-synthesis-anholonomy}，$[e_b,e_c]=C^d{}_{bc}\,e_d$，故
\begin{equation*}
 \dd\theta^a(e_b,e_c)
 =-\theta^a(C^d{}_{bc}\,e_d)
 =-C^d{}_{bc}\,\theta^a(e_d)
 =-C^d{}_{bc}\,\delta^a{}_d
 =-C^a{}_{bc}.
\end{equation*}
关键处用了 $\theta^a$ 的 $C^\infty$-线性把系数 $C^d{}_{bc}$ 提到外面，再用对偶性 $\theta^a(e_d)=\delta^a{}_d$ 收缩哑指标 $d$。

任意二形式都由其在基向量上的值唯一确定。因 $C^a{}_{bc}$ 对 $(b,c)$ 反对称，展开为楔积时需除以 $2$（重复计数修正）：
\begin{equation}\label{eq:math-dg-frame-structure-equation}
 \dd\theta^a
 =-\frac{1}{2}C^a{}_{bc}\,\theta^b\wedge\theta^c,
\end{equation}
即\eqnrefs{eq:math-dg-frame-structure-equation}。因子 $1/2$ 来自对反对称指标 $(b,c)$ 的重复计数：$\theta^b\wedge\theta^c$ 在 $(b,c)$ 和 $(c,b)$ 处取相反值，求和时每对出现两次。

典型例子。
\begin{itemize}
\item 完整标架 $e_a=\partial/\partial y^a$：偏导数交换，$C^a{}_{bc}\equiv0$，结构方程退化为 $\dd\theta^a=0$，即 $\dd(\dd y^a)=0$。
\item $\RR^2$ 上极坐标标架 $e_r=\partial_r$、$e_\theta=(1/r)\partial_\theta$，余标架 $\theta^r=\dd r$、$\theta^\theta=r\,\dd\theta$。由 $[e_r,e_\theta]=-(1/r)e_\theta$ 得 $C^\theta{}_{r\theta}=-C^\theta{}_{\theta r}=-1/r$，结构方程 $\dd\theta^\theta=\theta^r\wedge\theta^\theta$ 直接验证 $\dd(r\,\dd\theta)=\dd r\wedge\dd\theta$。
\item 左不变标架于 Lie 群 $G$：$C^a{}_{bc}$ 退化为 Lie 代数结构常数 $c^a{}_{bc}$，结构方程成为 Maurer--Cartan 方程 $\dd\theta^a=-\tfrac12 c^a{}_{bc}\theta^b\wedge\theta^c$。
\end{itemize}

非完整系数 $C^a{}_{bc}$ 在刚体力学中是旋转标架下的"假想力"（Coriolis、离心）系数；在广义相对论中，它是连接局部 Lorentz 标架与坐标系的运动学数据，与联络系数 $\omega^a{}_b$ 一起进入 Cartan 第一结构方程 $T^a=\dd\theta^a+\omega^a{}_b\wedge\theta^b$。
\end{derivation}

\eqnrefs{eq:math-dg-frame-structure-equation}并不是 Cartan 第一结构方程；它还没有引入联络和挠率。它只说明标架本身是否来自坐标。下一章加入联络以后，这个纯运动学恒等式会成为挠率结构方程的一部分。

在非坐标标架中，函数的外微分仍然没有任何额外修正：
\begin{equation*}\dd f=e_a(f)\theta^a.\end{equation*}
但是连续作两次 标架 导数 时不再交换：
\begin{equation*}[e_a,e_b]f
=e_a(e_bf)-e_b(e_af)
=C^c{}_{ab}e_cf.\end{equation*}
这已经表明 $C^c{}_{ab}$ 是“基向量作为微分算子”的非交换系数，而不是联络曲率。

若 $X=X^ae_a$、$Y=Y^ae_a$，Lie 括号的分量为
\begin{equation}
[X,Y]^c
=X^ae_a(Y^c)-Y^ae_a(X^c)
+C^c{}_{ab}X^aY^b.
\label{eq:lie-bracket-anholonomic-components}
\end{equation}
最后一项恰是从坐标基公式迁移到非完整标架时必须补上的结构项。

\begin{derivation}[非完整标架中的 Lie 括号 分量]
目标是把 Lie 括号 $[X,Y]$ 在非完整标架下的分量用结构函数 $C^c{}_{ab}$ 写出。先回忆对光滑函数 $f,g$ 与向量场 $X,Y$，Lie 括号作为微分算子作用在检验函数 $\varphi$ 上满足
\begin{equation*}
[fX,gY]\varphi
=fX(gY\varphi)-gY(fX\varphi).
\end{equation*}
由乘积法则
\begin{equation*}
X(gY\varphi)=X(g)\,Y\varphi+g\,X(Y\varphi),
\qquad
Y(fX\varphi)=Y(f)\,X\varphi+f\,Y(X\varphi),
\end{equation*}
代回得
\begin{align*}
[fX,gY]\varphi
&=fX(g)\,Y\varphi+fg\,X(Y\varphi)
 -gY(f)\,X\varphi-gf\,Y(X\varphi)\\
&=fg\bigl(X(Y\varphi)-Y(X\varphi)\bigr)
 +fX(g)\,Y\varphi-gY(f)\,X\varphi\\
&=fg[X,Y]\varphi+fX(g)\,Y\varphi-gY(f)\,X\varphi.
\end{align*}
由 $\varphi$ 任意，得到恒等式
\begin{equation*}
[fX,gY]=fg[X,Y]+fX(g)Y-gY(f)X.
\end{equation*}
于是令 $f=X^a$、$g=Y^b$、$X=e_a$、$Y=e_b$ 并对重复指标求和，
\begin{align*}
[X^ae_a,Y^be_b]
&=X^aY^b[e_a,e_b]
 +X^ae_a(Y^b)\,e_b
 -Y^be_b(X^a)\,e_a.
\end{align*}
由非完整标架定义 $[e_a,e_b]=C^c{}_{ab}e_c$，代回第一项，
\begin{equation*}
X^aY^b[e_a,e_b]
=X^aY^bC^c{}_{ab}e_c.
\end{equation*}
再把后两项的自由指标 $b$、$a$ 一致改写为 $c$，
\begin{align*}
X^ae_a(Y^b)\,e_b
&=X^ae_a(Y^c)\,e_c,\\
Y^be_b(X^a)\,e_a
&=Y^be_b(X^c)\,e_c.
\end{align*}
合并三项并提取公因子 $e_c$，
\begin{equation*}
[X,Y]
=\Bigl(
 X^ae_a(Y^c)-Y^be_b(X^c)
 +C^c{}_{ab}X^aY^b
\Bigr)e_c,
\end{equation*}
即\eqnrefs{eq:lie-bracket-anholonomic-components}。因此结构函数不仅出现在 $\dd\theta^a$ 中，也必然出现在所有以非坐标基表示的交换子计算中。
\end{derivation}

一组余标架是否来自局部坐标，可以直接由 外微分 判断。若存在局部函数 $y^a$ 使
\[
\theta^a=\dd y^a,
\]
则必有 $\dd\theta^a=0$。反过来，若每个 $\theta^a$ 都闭，则 Poincar\'e 引理局部给出 $\theta^a=\dd y^a$；由于余标架处处线性无关，Jacobian $\dd y^1\wedge\cdots\wedge\dd y^n$ 非零，故 $y^a$ 构成局部坐标。于是
\begin{equation*}C^a{}_{bc}=0\ \forall a,b,c
\quad\Longleftrightarrow\quad
\theta^a\text{ 局部为坐标余标架}.\end{equation*}
这是一组满秩标架的 Frobenius 条件；一般秩分布的版本将在后续章节系统处理。

\begin{exbox}[一个简单的非坐标标架]
在 $\RR^3$ 上取
\[
e_1=\partial_x+y\partial_z,
\qquad
e_2=\partial_y,
\qquad
e_3=\partial_z.
\]
直接计算
\[
[e_1,e_2]=-e_3,
\qquad
[e_1,e_3]=[e_2,e_3]=0.
\]
与它对偶的余标架为
\[
\theta^1=\dd x,
\qquad
\theta^2=\dd y,
\qquad
\theta^3=\dd z-y\,\dd x.
\]
于是
\[
\dd\theta^3
=-\dd y\wedge\dd x
=\theta^1\wedge\theta^2,
\]
恰好与 $C^3{}_{12}=-1$ 代入\eqnrefs{eq:math-dg-frame-structure-equation}得到的结果一致。
\end{exbox}

\section{标架变换下的结构函数}
结构函数不是张量分量，因为换标架时会出现 $\Lambda$ 的导数。把\eqnrefs{eq:math-dg-frame-change}代入 Lie 括号，得到
\[
\begin{aligned}
[e'_a,e'_b]
={}&[e_c\Lambda^c{}_a,e_d\Lambda^d{}_b]\\
={}&\Lambda^c{}_a\Lambda^d{}_b[e_c,e_d]
+\Lambda^c{}_a e_c(\Lambda^d{}_b)e_d
-\Lambda^d{}_b e_d(\Lambda^c{}_a)e_c.
\end{aligned}
\]
再用 $e_f=e'_r(\Lambda^{-1})^r{}_f$，可读出
\begin{align*}
C'^r{}_{ab}
=(\Lambda^{-1})^r{}_f\bigl(&\Lambda^c{}_a\Lambda^d{}_b C^f{}_{cd}
+\Lambda^c{}_a e_c(\Lambda^f{}_b)\\
&-\Lambda^c{}_b e_c(\Lambda^f{}_a)\bigr).
\end{align*}
第二行的导数项正是“结构函数不是张量”的原因。它描述的是所选基场自身怎样扭动，而不是一个独立于基的多线性映射。

若 $e^{(\alpha)}$ 与 $e^{(\beta)}$ 是两个局部标架，它们在重叠区满足
\[
e^{(\beta)}=e^{(\alpha)}\Lambda_{\alpha\beta}.
\]
在三重重叠上必须有
\begin{equation*}\Lambda_{\alpha\beta}\Lambda_{\beta\gamma}
=\Lambda_{\alpha\gamma}.\end{equation*}
所以换标架矩阵不是彼此独立的局部选择，而是切丛的 转移 上闭链。若所有 $\Lambda_{\alpha\beta}$ 都能同时写成
\[
\Lambda_{\alpha\beta}=h_\alpha^{-1}h_\beta,
\]
则可通过局部 规范 $h_\alpha$ 把 转移函数 全部消掉，得到全局标架；这等价于切丛 trivializable。不是每个流形都可平行化，因此“选一组全局 标架”本身可能受拓扑阻碍。

一个固定 张量 型 对应 $GL(n)$ 的一个表示。例如 $(1,1)$ 张量 在换标架下按 共轭 变换：
\[
T'^a{}_b
=(\Lambda^{-1})^a{}_cT^c{}_d\Lambda^d{}_b,
\]
而 陀螺 形式 的分量带 行列式 特征标：
\[
\theta'^1\wedge\cdots\wedge\theta'^n
=(\det\Lambda)^{-1}
\theta^1\wedge\cdots\wedge\theta^n.
\]
于是 取向 可以理解为把标架丛的结构群从 $GL(n,\RR)$ 约化到 $GL^+(n,\RR)$；Riemann 度量 则进一步把结构群约化到 $O(n)$，定向 Riemann 流形 约化到 $SO(n)$。后续 Spin 结构 正是在 $SO(n)$ 标架丛 上再寻求二重覆盖提升。

这一区分在后面十分重要：联络系数同样不是张量，也会在换标架时出现非齐次导数项；挠率与曲率却是张量。判断一个分量公式是否具有几何意义，不能只看指标形状，而要看它的换基律。

\section{非完整性与 Lie 群的原型}
若结构函数是常数，最重要的例子来自 Lie 群上的左不变标架。设 $G$ 为 Lie 群，$E_a$ 是 Lie 代数基，
\[
[E_b,E_c]=c^a{}_{bc}E_a.
\]
把 $E_a$ 左平移成全局左不变向量场，并记对偶左不变一形式为 $\vartheta^a$。由于结构系数不依赖位置，\eqnrefs{eq:math-dg-frame-structure-equation}化为 Maurer--Cartan 方程
\begin{equation}\label{eq:math-dg-maurer-cartan}
\boxed{
\dd\vartheta^a
=-\frac12c^a{}_{bc}\,\vartheta^b\wedge\vartheta^c.}
\end{equation}
它说明 Lie 代数的结构常数可以直接编码进一组一形式的外微分。
把 Maurer--Cartan 形式写成 Lie 代数 值一形式更紧凑。对矩阵 Lie 群 $G\subset GL(N)$，左 Maurer--Cartan 形式 为
\begin{equation*}\Theta=g^{-1}\dd g\in\Omega^1(G;\mathfrak g).\end{equation*}
直接使用
\[
\dd(g^{-1})=-g^{-1}(\dd g)g^{-1}
\]
得到
\begin{equation}
\dd\Theta+\Theta\wedge\Theta=0,
\label{eq:matrix-maurer-cartan-equation}
\end{equation}
其中矩阵乘法与 wedge 积 同时作用。若取 Lie 代数 基 $E_a$，写
\[
\Theta=\vartheta^aE_a,
\qquad
[E_b,E_c]=c^a{}_{bc}E_a,
\]
则\eqnrefs{eq:matrix-maurer-cartan-equation} 正好化为\eqnrefs{eq:math-dg-maurer-cartan}。

更重要的是，Maurer--Cartan 方程不仅是 Lie 群上已有形式的恒等式，也是一阶重构方程。设 $U$ 单连通，给定 $\mathfrak g$-值一形式 $A$ 满足
\begin{equation}
\dd A+A\wedge A=0.
\label{eq:mc-flat-integrability}
\end{equation}
则局部存在映射 $g:U\to G$ 使
\begin{equation}
A=g^{-1}\dd g.
\label{eq:mc-pure-gauge}
\end{equation}
在单连通区域上，适当全局条件下也可全局构造。方程\eqnrefs{eq:mc-flat-integrability} 是矩阵系统
\[
\dd g=gA
\]
路径无关可积的精确条件。

\begin{derivation}[Maurer--Cartan 方程是纯规范表示的可积条件，即\eqnrefs{eq:mc-flat-integrability} 对\eqnrefs{eq:mc-pure-gauge} 的充要性]
先证纯规范 $\Rightarrow$ 平坦。若 $A=g^{-1}\dd g$，由外微分的 Leibniz 规则 $\dd(g^{-1})=-g^{-1}(\dd g)g^{-1}$（对 $g\cdot g^{-1}=I$ 求微分），
\begin{align*}
 \dd A
 &=\dd(g^{-1})\wedge\dd g+g^{-1}\dd^2 g\\
 &=-g^{-1}(\dd g)g^{-1}\wedge\dd g+0\\
 &=-(g^{-1}\dd g)\wedge(g^{-1}\dd g)\\
 &=-A\wedge A,
\end{align*}
其中 $\dd^2g=0$（外微分的幂零性），第三个等号用了 $g^{-1}\dd g$ 的矩阵乘法与楔积的兼容性。因此 $\dd A+A\wedge A=0$，即\eqnrefs{eq:mc-flat-integrability} 必要。

再证平坦 $\Rightarrow$ 纯规范。固定基点 $x_0$ 与 $g_0\in G$，沿路径 $\gamma:[0,1]\to U$ 解矩阵 ODE
\begin{equation*}
 \frac{\dd g}{\dd t}
 =g(t)\,A_{\gamma(t)}(\dot\gamma(t)),
 \qquad
 g(0)=g_0.
\end{equation*}
这是一个线性 ODE，对每条路径有唯一解 $g_\gamma(1)\in G$。

若把路径作一个小矩形变形（先沿 $s$ 方向再沿 $t$ 方向，与先 $t$ 再 $s$），两种次序所得解的最低非平凡差由曲率
\begin{equation*}
 F_A=\dd A+A\wedge A
\end{equation*}
在矩形顶点处的值控制（显式小矩形计算见下文）。平坦条件 $F_A=0$ 使这个差为零，因此平行积分对可缩路径变形不敏感。单连通性进一步消除 同伦 障碍：任意两条具有相同端点的路径同伦，故 $g(x)$ 只依赖终点而不依赖所选路径，得到\eqnrefs{eq:mc-pure-gauge}。

取 $\RR^2$ 上小矩形 $[0,\epsilon]^2$，先 $s$ 后 $t$ 与先 $t$ 后 $s$ 的乘积相差
\begin{equation*}
 g_{\text{先}s\text{后}t}^{-1}g_{\text{先}t\text{后}s}
 =I+\epsilon^2 F_A(\partial_s,\partial_t)+O(\epsilon^3),
\end{equation*}
其中 $F_A$ 在矩形角点取值。这一展开来自 Taylor 展开 $A$ 与 BCH 公式保留 $\epsilon^2$ 阶。具体地，沿 $s$ 方向积分给 $g_s(\epsilon)=I+\epsilon A_s+\tfrac12\epsilon^2\partial_t A_s+O(\epsilon^3)$，沿 $t$ 方向给 $g_t(\epsilon)=I+\epsilon A_t+\tfrac12\epsilon^2\partial_s A_t+O(\epsilon^3)$，两种次序的乘积差含 $\epsilon^2(\partial_s A_t-\partial_t A_s+[A_s,A_t])=\epsilon^2 F_A(\partial_s,\partial_t)$。$F_A=0$ 时差异消失至 $O(\epsilon^3)$，把任意可缩大变形分解为小矩形即可归纳得路径无关。

若 $U$ 不单连通，平坦条件 $F_A=0$ 只保证局部纯规范，全局仍可能非平凡。基本群 $\pi_1(U)$ 通过 和乐 表示
\begin{equation*}
 \mathrm{Hol}_A:\pi_1(U,x_0)\to G,
 \qquad
 [\gamma]\mapsto g_\gamma(1)
\end{equation*}
给出全局信息。两个平坦联络 $A,A'$ 给同一 和乐 当且仅当差一个全局 规范 变换。因此平坦联络的同构类由 $\mathrm{Hom}(\pi_1(U),G)/G$ 参数化（模 空间）。这是规范理论中"真空模空间"的几何原型：即使场强为零，全局拓扑仍可区分不同规范真空。在 Chern--Simons 理论中，这一模空间正是三维流形上拓扑不变量的来源。

典型例子。
\begin{itemize}
\item $G=U(1)$：$A\in\Omega^1(U,\mathfrak u(1))\cong\Omega^1(U,\ii\RR)$，$F_A=\dd A$（因 Abel 群 $A\wedge A=0$）。平坦条件 $\dd A=0$ 即电磁场强度为零；局部 $A=\ii\,\dd\chi$ 即纯规范势，$\chi$ 为规范函数。这正是电磁学中"零场强 $\Rightarrow$ 局部纯规范"的几何表述。
\item $G=SU(2)$：$A=A_\mu^a\tau_a\,\dd x^\mu$（$\tau_a=\sigma_a/2\ii$），$F_A=\dd A+A\wedge A$ 含非 Abel 自相互作用项 $A\wedge A$。平坦 $SU(2)$ 联络局部可写 $A=g^{-1}\dd g$，$g:U\to SU(2)$；但全局可能存在非平凡 $SU(2)$ 丛（第二 Chern 类 $c_2\neq0$）即使局部平坦。
\item $G=\RR^n$（平移群）：$A$ 是普通闭一形式，平坦即 $\dd A=0$，纯 规范 $A=\dd f$ 对应势函数 $f$。这是 Hodge 分解中"闭形式 vs 恰当形式"的群论原型。
\end{itemize}

下一章中，非零曲率 $F_A\neq0$ 正是这个路径无关性失败的测量：沿不同路径的 和乐 产生非平凡群元素，即规范场有"磁通"。在物理上，Aharonov--Bohm 效应正是 $U(1)$ 平坦联络在非单连通区域（挖去磁通管）上具有非平凡 和乐 的体现：局部 $F=0$ 但全局 $A$ 不能写成纯规范，绕磁通管一圈积累相位 $e^{\ii\Phi/\hbar}$。
\end{derivation}

一般非坐标标架只是把常数 $c^a{}_{bc}$ 推广成位置依赖的 $C^a{}_{bc}(x)$。但任意给定的一组函数 $C^a{}_{bc}(x)$ 并不都能来自某个标架；它们必须满足 Lie 括号的 Jacobi 恒等式。这个微分约束也是 $\dd^2=0$ 在余标架结构方程上的具体表现。

\begin{derivation}[Jacobi 恒等式对结构函数的微分约束]
Lie 括号的 Jacobi 恒等式对结构函数 $C^c{}_{ab}$ 给出一组一阶微分约束；另一方面 $\dd^2=0$ 作用在结构方程上给出同一组约束。两条途径的等价性揭示"代数可积（Lie 代数）"与"微分可积（$\dd^2=0$）"是同一局部可积条件的两种语言。

由
\begin{equation*}
 [e_a,e_b]=C^c{}_{ab}e_c
\end{equation*}
和 Lie 括号的 Jacobi 恒等式
\begin{equation*}
 [e_a,[e_b,e_c]]+[e_b,[e_c,e_a]]+[e_c,[e_a,e_b]]=0,
\end{equation*}
第一项展开为（用 Lie 括号对第一因子 $C^\infty$-线性、对第二因子 Leibniz）
\begin{equation*}
\begin{aligned}
 [e_a,[e_b,e_c]]
 &=[e_a,C^r{}_{bc}e_r]\\
 &=e_a(C^r{}_{bc})e_r+C^r{}_{bc}[e_a,e_r]\\
 &=\bigl[e_a(C^d{}_{bc})+C^r{}_{bc}C^d{}_{ar}\bigr]e_d.
\end{aligned}
\end{equation*}
第二行用 Lie 括号对函数系数的 Leibniz 规则，第三行把 $[e_a,e_r]=C^d{}_{ar}e_d$ 代入并合并 $e_d$ 的系数。

对三个循环项 $(a,b,c)$、$(b,c,a)$、$(c,a,b)$ 求和，并利用 $\{e_d\}$ 线性无关得到
\begin{equation}\label{eq:math-dg-anholonomy-jacobi}
\boxed{
\sum_{\mathrm{cyc}(a,b,c)}
\left[
 e_a(C^d{}_{bc})+C^r{}_{bc}C^d{}_{ar}
\right]=0.}
\end{equation}
第一项 $e_a(C^d{}_{bc})$ 是结构函数的方向导数（微分部分），第二项 $C^r{}_{bc}C^d{}_{ar}$ 是结构函数的二次乘积（代数部分）。完整标架时 $C\equiv0$ 自动满足。

对\eqnrefs{eq:math-dg-frame-structure-equation}再作用一次 $\dd$，由 $\dd^2\equiv0$ 得
\begin{equation*}
 0=\dd^2\theta^d
 =-\frac12\dd C^d{}_{bc}\wedge\theta^b\wedge\theta^c
 -\frac12C^d{}_{bc}\dd(\theta^b\wedge\theta^c).
\end{equation*}
把 $\dd C^d{}_{bc}=e_a(C^d{}_{bc})\theta^a$（结构函数沿标架的方向导数）以及结构方程再次代入 $\dd\theta^b$，并读取 $\theta^a\wedge\theta^b\wedge\theta^c$ 的系数，恰好重新得到\eqnrefs{eq:math-dg-anholonomy-jacobi}。

取 $n=3$、$(a,b,c)=(1,2,3)$，Jacobi 约束为
\begin{equation*}
 e_1(C^d{}_{23})+e_2(C^d{}_{31})+e_3(C^d{}_{12})
 +C^r{}_{23}C^d{}_{1r}
 +C^r{}_{31}C^d{}_{2r}
 +C^r{}_{12}C^d{}_{3r}=0,
\end{equation*}
对 $d=1,2,3$ 各给一条方程。若 $C^d{}_{ab}$ 是常数（左不变标架于 Lie 群），微分项消失，只剩代数项
\begin{equation*}
 c^r{}_{bc}c^d{}_{ar}+c^r{}_{ca}c^d{}_{br}+c^r{}_{ab}c^d{}_{cr}=0,
\end{equation*}
即 Lie 代数结构常数的 Jacobi 恒等式。

这一可积条件是 Frobenius 定理与 Newlander--Nirenberg 定理的共同原型：一组分布可积当且仅当 Lie 括号封闭；一组近复结构可积当且仅当 Nijenhuis 张量消失。在规范理论中，Maurer--Cartan 方程 $\dd A+A\wedge A=0$ 是同一思想在主丛上的版本——"局部存在 $g$ 使 $A=g^{-1}\dd g$"的充要条件。
\end{derivation}


标架丛 $F(M)$ 的结构群天然是 $GL(n,\RR)$。许多“额外几何结构”都可以统一理解为把这个结构群约化到一个子群 $G\subset GL(n,\RR)$。所谓 $G$-结构，就是一个主 $G$-子丛
\begin{equation*}P_G\subset F(M)\end{equation*}
使得每个 纤维 中只保留满足指定几何条件的标架。

几个最重要的例子直接说明这一定义的统一性：
\begin{align*}
\text{orientation}
&\Longleftrightarrow GL(n,\RR)\to GL^+(n,\RR),\\
\text{Riemannian metric}
&\Longleftrightarrow GL(n,\RR)\to O(n),\\
\text{oriented Riemannian metric}
&\Longleftrightarrow GL(n,\RR)\to SO(n),\\
\text{almost complex structure }J
&\Longleftrightarrow GL(2m,\RR)\to GL(m,\CC),\\
\text{symplectic form }\omega
&\Longleftrightarrow GL(2m,\RR)\to Sp(2m,\RR).
\end{align*}
这里最后两行的含义分别是：允许的标架必须把固定模型复结构 $J_0$ 或模型 辛形式 $\omega_0$ 搬到当前点的 $J_p$、$\omega_p$。因此 张量场 与 结构-群 reduction 是同一几何信息的两种描述。

若 $G$ 是某个模型张量 $T_0$ 的稳定子群
\[
G=\{A\in GL(n,\RR):A\cdot T_0=T_0\},
\]
则一个 $G$-结构 等价于在 $M$ 上给定一个处处与 $T_0$ 属于同一 $GL(n)$-轨道 的 张量场 $T$。例如 Euclidean 内积的稳定子群是 $O(n)$，所以 Riemann 度量 正是 $O(n)$-结构；标准复结构的稳定子群是 $GL(m,\CC)$，所以 殆复结构 是 $GL(m,\CC)$-结构。

但“约化结构群”和“结构可积”必须严格区分。一个 殆复结构 只要求
\[
J^2=-I,
\]
并不自动来自复坐标。其 integrability 障碍 是 Nijenhuis 张量
\begin{equation*}N_J(X,Y)
=[JX,JY]-J[JX,Y]-J[X,JY]-[X,Y].\end{equation*}
Newlander--Nirenberg 定理说明
\[
N_J=0
\quad\Longleftrightarrow\quad
J\text{ 局部来自复坐标图册}.
\]
这与下一章的 挠率、再下一章的 Frobenius 障碍 具有相同哲学：给定 pointwise 代数 结构 只是第一步，真正的局部几何还要检查其一阶微分可积条件。

同样地，一个非退化二形式 $\omega$ 把结构群约化到 $Sp(2m,\RR)$，但只有再满足
\begin{equation*}\dd\omega=0\end{equation*}
时才是 辛 结构。Darboux 定理 随后说明所有 辛形式 局部都可以化为
\[
\omega=\sum_{a=1}^m\dd q^a\wedge\dd p_a.
\]
因此 辛 几何没有类似 Riemann 曲率那样的局部零阶不变量；其局部非平凡性首先受 closedness 这一一阶微分条件控制。

G-结构 还解释了“兼容联络”意味着什么。若 $\omega$ 是 $F(M)$ 上的 $\mathfrak{gl}(n)$-值联络一形式，一个联络保持 $G$-结构，当且仅当限制到 $P_G$ 后它取值于 Lie 代数 $\mathfrak g$：
\begin{equation*}\omega|_{P_G}\in\Omega^1(P_G;\mathfrak g).\end{equation*}
Riemann 情形 $G=O(n)$ 时，这就是联络形式反对称；若再要求 挠率-自由，便唯一得到 Levi--Civita 联络。其他 $G$-结构 未必存在唯一 挠率-自由 compatible 联络；剩余的不可消去 挠率 成分称为 intrinsic 挠率，它刻画该 $G$-结构 偏离“最理想局部模型”的一阶程度。

这套语言把后续多种结构统一在一张图里：
\[
\text{frame bundle}
\longrightarrow
\text{structure-group reduction}
\longrightarrow
\text{compatible connection}
\longrightarrow
\text{torsion/curvature/holonomy}.
\]
因此标架并非只为写 Cartan 公式服务，它实际上是联结 张量 几何、规范 几何 与特殊 和乐 的共同母空间。

这一章至此仍未回答“如何把一个向量沿曲线平行搬运”以及“怎样比较邻点上的向量”。结构函数只刻画基场的非完整性。下一章引入 仿射 联络 后，平行移动、挠率与曲率才成为新的几何数据。
